\documentclass[10pt,a4paper]{amsart}
\usepackage{amsmath,amssymb,amsthm,graphicx}
% For commutative diagrams we use Xy-pic
\usepackage[all]{xy}

% We use 2cell for 2-commutative diagrams.
\xyoption{2cell}
% Theorem environments.
%
\theoremstyle{plain}
\newtheorem{theorem}[subsection]{Theorem}
\newtheorem{proposition}[subsection]{Proposition}
\newtheorem{lemma}[subsection]{Lemma}

\theoremstyle{definition}
\newtheorem{definition}[subsection]{Definition}
\newtheorem{example}[subsection]{Example}
\newtheorem{exercise}[subsection]{Exercise}
\newtheorem{situation}[subsection]{Situation}

\theoremstyle{remark}
\newtheorem{remark}[subsection]{Remark}
\newtheorem{remarks}[subsection]{Remarks}

\numberwithin{equation}{subsection}

% Macros
%
\def\lim{\mathop{\mathrm{lim}}\nolimits}
\def\colim{\mathop{\mathrm{colim}}\nolimits}
\def\Spec{\mathop{\mathrm{Spec}}}
\def\Hom{\mathop{\mathrm{Hom}}\nolimits}
\def\Ext{\mathop{\mathrm{Ext}}\nolimits}
\def\SheafHom{\mathop{\mathcal{H}\!\mathit{om}}\nolimits}
\def\SheafExt{\mathop{\mathcal{E}\!\mathit{xt}}\nolimits}
\def\Sch{\mathit{Sch}}
\def\Mor{\mathop{\mathrm{Mor}}\nolimits}
\def\Ob{\mathop{\mathrm{Ob}}\nolimits}
\def\Sh{\mathop{\mathit{Sh}}\nolimits}
\def\NL{\mathop{N\!L}\nolimits}
\def\CH{\mathop{\mathrm{CH}}\nolimits}
\def\proetale{{pro\text{-}\acute{e}tale}}
\def\etale{{\acute{e}tale}}
\def\QCoh{\mathit{QCoh}}
\def\Ker{\mathop{\mathrm{Ker}}}
\def\Im{\mathop{\mathrm{Im}}}
\def\Coker{\mathop{\mathrm{Coker}}}
\def\Coim{\mathop{\mathrm{Coim}}}

% Boxtimes
%
\DeclareMathSymbol{\boxtimes}{\mathbin}{AMSa}{"02}

%
% Macros for moduli stacks/spaces
%
\def\QCohstack{\mathcal{QC}\!\mathit{oh}}
\def\Cohstack{\mathcal{C}\!\mathit{oh}}
\def\Spacesstack{\mathcal{S}\!\mathit{paces}}
\def\Quotfunctor{\mathrm{Quot}}
\def\Hilbfunctor{\mathrm{Hilb}}
\def\Curvesstack{\mathcal{C}\!\mathit{urves}}
\def\Polarizedstack{\mathcal{P}\!\mathit{olarized}}
\def\Complexesstack{\mathcal{C}\!\mathit{omplexes}}
% \Pic is the operator that assigns to X its picard group, usage \Pic(X)
% \Picardstack_{X/B} denotes the Picard stack of X over B
% \Picardfunctor_{X/B} denotes the Picard functor of X over B
\def\Pic{\mathop{\mathrm{Pic}}\nolimits}
\def\Picardstack{\mathcal{P}\!\mathit{ic}}
\def\Picardfunctor{\mathrm{Pic}}
\def\Deformationcategory{\mathcal{D}\!\mathit{ef}}
\usepackage{cleveref}
\crefname{originalreference}{original reference}{original references}
\Crefname{originalreference}{Original reference}{Original references}
\makeatletter
\newlabel{lemma-base-change-q-injective}{{0F03}{source}{}{sourcepointer-1}{}}
\newlabel{lemma-base-change-q-injective@cref}{{[originalreference][0][]0F03}{[1][0][]source}{}{}{}}
\makeatother
\crefname{originalreference}{original reference}{original references}
\Crefname{originalreference}{Original reference}{Original references}
\makeatletter
\newlabel{lemma-integral-pushforward-commutes-with-base-change}{{0EYP}{source}{}{sourcepointer-2}{}}
\newlabel{lemma-integral-pushforward-commutes-with-base-change@cref}{{[originalreference][0][]0EYP}{[1][0][]source}{}{}{}}
\makeatother
\crefname{originalreference}{original reference}{original references}
\Crefname{originalreference}{Original reference}{Original references}
\makeatletter
\newlabel{lemma-what-integral}{{04C2}{source}{}{sourcepointer-3}{}}
\newlabel{lemma-what-integral@cref}{{[originalreference][0][]04C2}{[1][0][]source}{}{}{}}
\makeatother
\crefname{originalreference}{original reference}{original references}
\Crefname{originalreference}{Original reference}{Original references}
\makeatletter
\newlabel{remark-base-change-holds}{{0F02}{source}{}{sourcepointer-4}{}}
\newlabel{remark-base-change-holds@cref}{{[originalreference][0][]0F02}{[1][0][]source}{}{}{}}
\makeatother
\crefname{originallemma}{original lemma}{original lemmas}
\Crefname{originallemma}{Original lemma}{Original lemmas}
\makeatletter
\newlabel{sites-cohomology-lemma-relative-Leray}{{0734}{source}{}{sourcepointer-5}{}}
\newlabel{sites-cohomology-lemma-relative-Leray@cref}{{[originallemma][0][]0734}{[1][0][]source}{}{}{}}
\makeatother
\title{Base change for a higher direct image can be checked after an integral surjective base change}
\author{The Stacks Project Authors}
\date{stacks-project:etale-cohomology.tex:local-snapshot-20261011; GNU FDL 1.2 or later}
\begin{document}
\maketitle
\noindent\textit{Editorial note.} All delivered mathematical text is the author's own, copied from the pinned Stacks Project chapter file etale-cohomology.tex; no statement or proof step was rewritten, reordered or removed.\par\medskip
\noindent\textit{Editorial note.} The counted proof is the author's complete argument: the injectivity of the adjunction map along the surjective integral morphism, the reduction to the pushed forward coefficient sheaf, and the three identifications coming from the compatibility with integral base change and the relative Leray spectral sequence.\par\medskip
\noindent\textit{Editorial note.} The label lemma-base-change-q-integral-top names the counted statement, and the proof environment immediately following it in the source is delivered verbatim as its proof.\par\medskip
\noindent\textit{Editorial note.} The counted statement is carried by the author's own lemma environment in the chapter source.\par\medskip
\noindent\textit{Editorial note.} Pre-open gate: python operations/agent-authoring/tools/precard\_gate.py --source etale-cohomology.tex --label lemma-base-change-q-integral-top --start 16938 --end 16994 --local-id 3568 returned RESULT=OK (L3568 free, occupancy FREE, public interval clear). Candidate taken from lane-s4/halfB-priority.json.\par\medskip
\noindent\textit{Editorial note.} The statement and proof are copied from the pinned chapter file; the cited inputs are kept as external pointers with their canonical Stacks tags.\par\medskip
\noindent\textit{Allowed standard prerequisites (editorial).}\begin{itemize}
\item Injectivity of the adjunction map along a surjective morphism. Same chapter, the standard fact that the adjunction map F -> pi\_*pi\textasciicircum{}\{-1\}F is injective when pi is surjective, checked on stalks. Supplies the injectivity of F -> pi\_*pi\textasciicircum{}\{-1\}F' used to reduce the assertion to the pullback of that pushforward.
\item Base change for an injective coefficient sheaf. Same chapter, the lemma reducing the base change property for a coefficient sheaf to the corresponding property for a subsheaf. Supplies the reduction of the base change map for F to the base change map for pi\_*pi\textasciicircum{}\{-1\}F'.
\item Commutation of the higher direct images with an integral base change, and the relative Leray spectral sequence. Same chapter, the lemmas that the pushforward along an integral morphism commutes with base change and that the higher direct images along a composite behave as in the relative Leray spectral sequence. Supply the three identifications R\textasciicircum{}q g\_* pi\_* F' = R\textasciicircum{}q g'\_* F', e\textasciicircum{}\{-1\} pi\_* F' = pi'\_*(e')\textasciicircum{}\{-1\}F' and R\textasciicircum{}q h\_* pi'\_*(e')\textasciicircum{}\{-1\}F' = R\textasciicircum{}q h'\_*(e')\textasciicircum{}\{-1\}F'.
\end{itemize}\medskip
\ifdefined\hypertarget\hypertarget{sourcepointer-1}{}\fi
\noindent\textit{Source reference.} Original reference 0F03: sources/stacks-etale-cohomology/pinned-original-source.tex, source lines 16869--16869. Prerequisite pointer invoked at proof line(s) 16979; it names a standard fact of the same chapter that the counted argument uses, and every step of the counted proof built on it is retained verbatim in the delivered body.\par\medskip
\ifdefined\hypertarget\hypertarget{sourcepointer-2}{}\fi
\noindent\textit{Source reference.} Original reference 0EYP: sources/stacks-etale-cohomology/pinned-original-source.tex, source lines 8350--8350. Prerequisite pointer invoked at proof line(s) 16991; it names a standard fact of the same chapter that the counted argument uses, and every step of the counted proof built on it is retained verbatim in the delivered body.\par\medskip
\ifdefined\hypertarget\hypertarget{sourcepointer-3}{}\fi
\noindent\textit{Source reference.} Original reference 04C2: sources/stacks-etale-cohomology/pinned-original-source.tex, source lines 6056--6056. Prerequisite pointer invoked at proof line(s) 16992; it names a standard fact of the same chapter that the counted argument uses, and every step of the counted proof built on it is retained verbatim in the delivered body.\par\medskip
\ifdefined\hypertarget\hypertarget{sourcepointer-4}{}\fi
\noindent\textit{Source reference.} Original reference 0F02: sources/stacks-etale-cohomology/pinned-original-source.tex, source lines 16848--16848. Prerequisite pointer invoked at proof line(s) 16940; it names a standard fact of the same chapter that the counted argument uses, and every step of the counted proof built on it is retained verbatim in the delivered body.\par\medskip
\ifdefined\hypertarget\hypertarget{sourcepointer-5}{}\fi
\noindent\textit{Source reference.} Original lemma 0734: sources/stacks-chow/tags.txt, source lines 9048--9048. Prerequisite pointer invoked at proof line(s) 16993; it names a standard fact of a sibling chapter that the counted argument uses, and every step of the counted proof built on it is retained verbatim in the delivered body.\par\medskip
\begin{lemma}
\label{lemma-base-change-q-integral-top}
With $f : X \to S$ and $n$ as in Remark \ref{remark-base-change-holds}
assume for some $q \geq 1$ we have $BC(f, n, q - 1)$. Consider
commutative diagrams
$$
\vcenter{
\xymatrix{
X \ar[d]_f &
X' \ar[d]_{f'} \ar[l] &
Y \ar[l]^h \ar[d]^e &
Y' \ar[l]^{\pi'} \ar[d]^{e'} \\
S &
S' \ar[l] &
T \ar[l]_g &
T' \ar[l]_\pi
}
}
\quad\text{and}\quad
\vcenter{
\xymatrix{
X' \ar[d]_{f'} & & Y' \ar[ll]^{h' = h \circ \pi'} \ar[d]^{e'} \\
S' & & T' \ar[ll]_{g' = g \circ \pi}
}
}
$$
where all squares are cartesian, $g$ quasi-compact and quasi-separated, and
$\pi$ is integral surjective. Let $\mathcal{F}$ be an abelian sheaf
on $T_\etale$ annihilated by $n$ and set $\mathcal{F}' = \pi^{-1}\mathcal{F}$.
If the base change map
$$
(f')^{-1}R^qg'_*\mathcal{F}' \longrightarrow R^qh'_*(e')^{-1}\mathcal{F}'
$$
is an isomorphism, then the base change map
$(f')^{-1}R^qg_*\mathcal{F} \to R^qh_*e^{-1}\mathcal{F}$
is an isomorphism.
\end{lemma}

\begin{proof}
Observe that $\mathcal{F} \to \pi_*\pi^{-1}\mathcal{F}'$ is injective
as $\pi$ is surjective (check on stalks). Thus by
Lemma \ref{lemma-base-change-q-injective}
we see that it suffices to show that the base change map
$$
(f')^{-1}R^qg_*\pi_*\mathcal{F}'
\longrightarrow
R^qh_*e^{-1}\pi_*\mathcal{F}'
$$
is an isomorphism. This follows from the assumption because
we have $R^qg_*\pi_*\mathcal{F}' = R^qg'_*\mathcal{F}'$,
we have $e^{-1}\pi_*\mathcal{F}' =\pi'_*(e')^{-1}\mathcal{F}'$, and
we have $R^qh_*\pi'_*(e')^{-1}\mathcal{F}' = R^qh'_*(e')^{-1}\mathcal{F}'$.
This follows from Lemmas
\ref{lemma-integral-pushforward-commutes-with-base-change} and
\ref{lemma-what-integral} and the relative leray spectral sequence
(Cohomology on Sites, Lemma \ref{sites-cohomology-lemma-relative-Leray}).
\end{proof}

\end{document}
